% GENERATED FILE. Edit canonical lessons, then run npm run generate:books.
% canonical-source-hash: fnv1a64:0f4e277888f4d856
% canonical-lessons: SA-C492-garbhini, SA-C492-prasute, SA-C492-janma, SA-C492-navajatah, SA-C492-dola

\chapter{A New Baby}
\label{ch:sa-a-new-baby}
\hlchaptermodality{492}

\begin{chapteropening}
\textbf{By the end of this chapter:} \emph{I can talk about the birth of a baby.}

Complete the last lesson of chapter 492: I can talk about the birth of a baby.
\end{chapteropening}

\section[\texorpdfstring{\sk{गर्भिणी} (garbhiṇī)}{गर्भिणी (garbhiṇī)}]{\sk{गर्भिणी} (garbhiṇī) — a pregnant woman}

\label{lesson:SA-C492-garbhini}



\begin{quote}
Before the new one: say the Sanskrit for funeral rites, then the Sanskrit for a rite for the ancestors.
\end{quote}

\subsection*{You'll want to know: \sk{गर्भिणी}}
\textbf{\sk{गर्भिणी}} — \emph{garbhiṇī} — \textquotedblleft{}a pregnant woman\textquotedblright{}.

\begin{quote}
\sk{गर्भिणी} \sk{पथ्यं} \sk{भोजनं} \sk{खादति।}
\end{quote}

\begin{quote}
\emph{garbhiṇī pathyaṃ bhojanaṃ khādati.}
\end{quote}

\begin{quote}
The pregnant woman eats wholesome food.
\end{quote}

A feminine noun.

\begin{cousinweb}
\textbf{\sk{गर्भिणी}} (\emph{garbhiṇī}) is built on \emph{garbha}, \textquotedblleft{}the womb\textquotedblright{}; MW's \emph{garbhin} is \textquotedblleft{}pregnant\textquotedblright{}, and this is its feminine (MW).
\end{cousinweb}

\begin{grammarlens}[title={What you've built}]
One more word for describing what you lived and felt.
\end{grammarlens}

\subsection*{Your turn}
\begin{itemize}
\raggedright
  \item \emph{Say it:} \emph{garbhiṇī}
  \item \emph{Say it:} \emph{garbhiṇī}, once more
  \item \emph{Say it:} \emph{garbhiṇī pathyaṃ bhojanaṃ khādati.} — The pregnant woman eats wholesome food.
  \item \emph{Recall:} say the Sanskrit for funeral rites, then the Sanskrit for a rite for the ancestors, then say \emph{garbhiṇī} again
\end{itemize}

\begin{tcolorbox}[breakable,colback=teal!4,colframe=teal!35!black,title={Before you move on}]
What does \sk{गर्भिणी} mean? (\textquotedblleft{}A pregnant woman\textquotedblright{}.) Say it once more.
\end{tcolorbox}
\section[\texorpdfstring{\sk{प्रसूते} (prasūte)}{प्रसूते (prasūte)}]{\sk{प्रसूते} (prasūte) — gives birth}

\label{lesson:SA-C492-prasute}



\begin{quote}
Before the new one: say the Sanskrit for a rite for the ancestors, then the Sanskrit for a pregnant woman.
\end{quote}

\subsection*{You'll want to know: \sk{प्रसूते}}
\textbf{\sk{प्रसूते}} — \emph{prasūte} — \textquotedblleft{}gives birth\textquotedblright{}.

\begin{quote}
\sk{धेनुः} \sk{वत्सं} \sk{प्रसूते।}
\end{quote}

\begin{quote}
\emph{dhenuḥ vatsaṃ prasūte.}
\end{quote}

\begin{quote}
The cow gives birth to a calf.
\end{quote}

Past (she …): \sk{सा} \sk{पुत्रं} \sk{प्रसूतवती}, \emph{sā putraṃ prasūtavatī}.

\begin{cousinweb}
\textbf{\sk{प्रसूते}} (\emph{prasūte}) is built from \emph{pra} and the root √sū, \textquotedblleft{}to beget, to bear\textquotedblright{}; MW gives \textquotedblleft{}to procreate, bring forth\textquotedblright{} (MW).
\end{cousinweb}

\begin{grammarlens}[title={What you've built}]
One more word for describing what you lived and felt.
\end{grammarlens}

\subsection*{Your turn}
\begin{itemize}
\raggedright
  \item \emph{Say it:} \emph{prasūte}
  \item \emph{Say it:} \emph{prasūte}, once more
  \item \emph{Say it:} \emph{dhenuḥ vatsaṃ prasūte.} — The cow gives birth to a calf.
  \item \emph{Recall:} say the Sanskrit for a rite for the ancestors, then the Sanskrit for a pregnant woman, then say \emph{prasūte} again
\end{itemize}

\begin{tcolorbox}[breakable,colback=teal!4,colframe=teal!35!black,title={Before you move on}]
What does \sk{प्रसूते} mean? (\textquotedblleft{}Gives birth\textquotedblright{}.) Say it once more.
\end{tcolorbox}
\section[\texorpdfstring{\sk{जन्म} (janma)}{जन्म (janma)}]{\sk{जन्म} (janma) — birth}

\label{lesson:SA-C492-janma}



\begin{quote}
Before the new one: say the Sanskrit for a pregnant woman, then the Sanskrit for gives birth.
\end{quote}

\subsection*{You'll want to know: \sk{जन्म}}
\textbf{\sk{जन्म}} — \emph{janma} — \textquotedblleft{}birth\textquotedblright{}.

\begin{quote}
\sk{मम} \sk{भ्रातुः} \sk{जन्म} \sk{रात्रौ} \sk{अभवत्।}
\end{quote}

\begin{quote}
\emph{mama bhrātuḥ janma rātrau abhavat.}
\end{quote}

\begin{quote}
My brother was born at night.
\end{quote}

A neuter noun.

\begin{cousinweb}
\textbf{\sk{जन्म}} (\emph{janma}) comes from the root √jan, \textquotedblleft{}to be born\textquotedblright{}, the root of \textbf{\sk{जायते}}; MW gives \textquotedblleft{}birth, production, origin\textquotedblright{} (MW).
\end{cousinweb}

\begin{grammarlens}[title={What you've built}]
One more word for describing what you lived and felt.
\end{grammarlens}

\subsection*{Your turn}
\begin{itemize}
\raggedright
  \item \emph{Say it:} \emph{janma}
  \item \emph{Say it:} \emph{janma}, once more
  \item \emph{Say it:} \emph{mama bhrātuḥ janma rātrau abhavat.} — My brother was born at night.
  \item \emph{Recall:} say the Sanskrit for a pregnant woman, then the Sanskrit for gives birth, then say \emph{janma} again
\end{itemize}

\begin{tcolorbox}[breakable,colback=teal!4,colframe=teal!35!black,title={Before you move on}]
What does \sk{जन्म} mean? (\textquotedblleft{}Birth\textquotedblright{}.) Say it once more.
\end{tcolorbox}
\section[\texorpdfstring{\sk{नवजातः} (navajātaḥ)}{नवजातः (navajātaḥ)}]{\sk{नवजातः} (navajātaḥ) — newborn}

\label{lesson:SA-C492-navajatah}



\begin{quote}
Before the new one: say the Sanskrit for gives birth, then the Sanskrit for birth.
\end{quote}

\subsection*{You'll want to know: \sk{नवजातः}}
\textbf{\sk{नवजातः}} — \emph{navajātaḥ} — \textquotedblleft{}newborn\textquotedblright{}.

\begin{quote}
\sk{नवजातः} \sk{शिशुः} \sk{बहु} \sk{स्वपिति।}
\end{quote}

\begin{quote}
\emph{navajātaḥ śiśuḥ bahu svapiti.}
\end{quote}

\begin{quote}
A newborn baby sleeps a lot.
\end{quote}

\begin{cousinweb}
\textbf{\sk{नवजातः}} (\emph{navajātaḥ}) is built from \textbf{\sk{नव}}, \textquotedblleft{}new\textquotedblright{}, and \emph{jāta}, \textquotedblleft{}born\textquotedblright{}, from the root √jan; MW gives \textquotedblleft{}recently born, new\textquotedblright{} (MW).
\end{cousinweb}

\begin{grammarlens}[title={What you've built}]
One more word for describing what you lived and felt.
\end{grammarlens}

\subsection*{Your turn}
\begin{itemize}
\raggedright
  \item \emph{Say it:} \emph{navajātaḥ}
  \item \emph{Say it:} \emph{navajātaḥ}, once more
  \item \emph{Say it:} \emph{navajātaḥ śiśuḥ bahu svapiti.} — A newborn baby sleeps a lot.
  \item \emph{Recall:} say the Sanskrit for gives birth, then the Sanskrit for birth, then say \emph{navajātaḥ} again
\end{itemize}

\begin{tcolorbox}[breakable,colback=teal!4,colframe=teal!35!black,title={Before you move on}]
What does \sk{नवजातः} mean? (\textquotedblleft{}Newborn\textquotedblright{}.) Say it once more.
\end{tcolorbox}
\section[\texorpdfstring{\sk{दोला} (dolā)}{दोला (dolā)}]{\sk{दोला} (dolā) — a swing; a cradle}

\label{lesson:SA-C492-dola}



\begin{quote}
Before the new one: say the Sanskrit for birth, then the Sanskrit for newborn.
\end{quote}

\subsection*{You'll want to know: \sk{दोला}}
\textbf{\sk{दोला}} — \emph{dolā} — \textquotedblleft{}a swing; a cradle\textquotedblright{}.

\begin{quote}
\sk{शिशुः} \sk{दोलायां} \sk{शेते।}
\end{quote}

\begin{quote}
\emph{śiśuḥ dolāyāṃ śete.}
\end{quote}

\begin{quote}
The baby lies in the cradle.
\end{quote}

A feminine noun.

\begin{cousinweb}
\textbf{\sk{दोला}} (\emph{dolā}) MW gives \textquotedblleft{}litter, hammock, palanquin, swing\textquotedblright{}; MW gives no derivation, so its origin is uncertain (MW).
\end{cousinweb}

\begin{grammarlens}[title={What you've built}]
One more word for describing what you lived and felt.
\end{grammarlens}

\subsection*{Your turn}
\begin{itemize}
\raggedright
  \item \emph{Say it:} \emph{dolā}
  \item \emph{Say it:} \emph{dolā}, once more
  \item \emph{Say it:} \emph{śiśuḥ dolāyāṃ śete.} — The baby lies in the cradle.
  \item \emph{Recall:} say the Sanskrit for birth, then the Sanskrit for newborn, then say \emph{dolā} again
\end{itemize}

\begin{tcolorbox}[breakable,colback=teal!4,colframe=teal!35!black,title={Before you move on}]
What does \sk{दोला} mean? (\textquotedblleft{}A swing; a cradle\textquotedblright{}.) Say it once more.
\end{tcolorbox}
